\documentclass[10pt,a4paper]{amsart}
\usepackage[margin=19mm]{geometry}
\usepackage{amsmath,amssymb,mathtools}
\usepackage[scr=rsfs,cal=euler]{mathalfa}
\usepackage{xfrac}
\usepackage[unicode,hidelinks,bookmarks=false]{hyperref}
\newcommand{\ie}{i.e.\ }
\newcommand{\cf}{cf.\ }
\newcommand{\resp}{respectively\ }
\newcommand{\Subsection}{\subsection}
\providecommand{\nobreakdash}{\nobreak\hbox{-}}
\setlength{\emergencystretch}{2em}
\multlinegap0pt
\usepackage{amssymb}
\usepackage{xcolor}
\usepackage{tikz}
\usepackage{tcolorbox}
\newtheorem{theorem}{Theorem}
\title{Presentations and Representations of the Multi-Virtual Twin Group and Associated Subgroups}
\author{Vaibhav Keshari, Taher I. Mayassi, Madeti Prabhakar, Mohamad N. Nasser}
\date{Source: arXiv:2605.13090v1 (CC BY 4.0)}
\begin{document}
\maketitle

\noindent\emph{Source and licence.} Presentations and Representations of the Multi-Virtual Twin Group and Associated Subgroups. Version arXiv:2605.13090v1 (CC BY 4.0), distributed by arXiv under the Creative Commons Attribution 4.0 International licence, http://creativecommons.org/licenses/by/4.0/, as recorded in the arXiv licence metadata for this exact version and shown on the abstract page https://arxiv.org/abs/2605.13090v1. Full licence notice: \texttt{licenses/arxiv-2605.13090-CC-BY-4.0.txt}.\par\smallskip

\noindent\textit{Editorial scope.} Counted unit: the article's theorem theorem (source0.tex lines 1150--1167). The article's complete original proof of it is delivered with it (source0.tex lines 1169--1381, one \texttt{proof} environment of 213 lines). No mathematics is written, repaired or re-derived: the statement and the proof are the article's own lines, in the article's own notation, and no retained line is substituted or re-flowed.  Retained uncounted as the source-local context the proof needs: lines 651--931.  Nothing was omitted from any retained span, so the package carries no omission record.  No retained line contains a citation, so the wrapper carries no bibliography and no bibliography entry.  No retained line contains a figure environment, an image-inclusion command or drawing code, so no image is delivered and the package declares no asset dependency.  Editorial formatting only, in addition: an AMS \texttt{amsart} preamble that restates the article's own declarations and macros used by the retained lines, namely the environments \texttt{theorem} on separate counters, together with the standard packages those lines need.  The retained environments print in this document as Theorem 1 in order of appearance, whereas the article prints the same items with the numbering of its own full document.  The complete native source of the article as distributed in the arXiv e-print for this exact version is bundled unchanged as \texttt{sources/200477/source0.tex}.
\begin{theorem} \label{thm 41}
    Let $\zeta: M_kVT_n \longrightarrow \mathrm{GL}_{n}(\mathbb{C})$ be a complex homogeneous 2-local representation of $M_kVT_n$ for  \(k \geq 1\) and \(n \geq 3\). Then, $\zeta$ is equivalent to one of the following eight representations.
    \begin{itemize}
    \item[(1)] $\zeta_1: M_kVT_n \longrightarrow \mathrm{GL}_n(\mathbb{C})$ such that
$$\zeta_1(s_i)=\zeta_1(\rho_i^{\alpha})= I_n$$
    for $1\leq i \leq n-1, 0 \leq \alpha \leq k-1.$ 
    %\vspace{0.2cm}
%%%%%%%%%%%%%%%%%%%%%%%%%%%%%%%%%%%%%%%%%%%%%%%%%%%%%%%%%%%%%%%%%%%%%%%%%%%
\item[(2)] $\zeta_2: M_kVT_n \longrightarrow \mathrm{GL}_n(\mathbb{C})$ such that
{\small
$$\zeta_2(s_i)=\left( \begin{array}{c|@{}c|c@{}}
   \begin{matrix}
     I_{i-1} 
   \end{matrix} 
      & \textbf{0} & \textbf{0} \\
      \hline
    \textbf{0} &\hspace{0.2cm} \begin{matrix}
   		1 & 0\\
   		0 & 1\\
   		\end{matrix}  & \textbf{0}  \\
\hline
\textbf{0} & \textbf{0} & I_{n-i-1}
\end{array} \right) \text{ and } \zeta_2(\rho_i^{\alpha})=\left( \begin{array}{c|@{}c|c@{}}
   \begin{matrix}
     I_{i-1} 
   \end{matrix} 
      & \textbf{0} & \textbf{0} \\
      \hline
    \textbf{0} &\hspace{0.2cm} \begin{matrix}
   		0 & \dfrac{1}{y_{\alpha}}\\
   		y_{\alpha} & 0\\
   		\end{matrix}  & \textbf{0}  \\
\hline
\textbf{0} & \textbf{0} & I_{n-i-1}
\end{array} \right)$$
}
for $1\leq i \leq n-1, 0 \leq \alpha \leq k-1, \text{ where } y_{\alpha}\in \mathbb{C}^*.$
%\vspace{0.2cm}
%%%%%%%%%%%%%%%%%%%%%%%%%%%%%%%%%%%%%%%%%%%%%%%%%%%%%%%%%%%%%%%%%%%%%%%%%%%
\item[(3)] $\zeta_3: M_kVT_n \longrightarrow \mathrm{GL}_n(\mathbb{C})$ such that
{\small
$$\zeta_3(s_i)=\left( \begin{array}{c|@{}c|c@{}}
   \begin{matrix}
     I_{i-1} 
   \end{matrix} 
      & \textbf{0} & \textbf{0} \\
      \hline
    \textbf{0} &\hspace{0.2cm} \begin{matrix}
   		1 & 0\\
   		0 & -1\\
   		\end{matrix}  & \textbf{0}  \\
\hline
\textbf{0} & \textbf{0} & I_{n-i-1}
\end{array} \right) \text{ and } \zeta_3(\rho_i^{\alpha})=\left( \begin{array}{c|@{}c|c@{}}
   \begin{matrix}
     I_{i-1} 
   \end{matrix} 
      & \textbf{0} & \textbf{0} \\
      \hline
    \textbf{0} &\hspace{0.2cm} \begin{matrix}
   		0 & \dfrac{1}{y_{\alpha}}\\
   		y_{\alpha} & 0\\
   		\end{matrix}  & \textbf{0}  \\
\hline
\textbf{0} & \textbf{0} & I_{n-i-1}
\end{array} \right)$$
}
$\text{ for } 1\leq i \leq n-1, 0 \leq \alpha \leq k-1, \text{ where } y_{\alpha}\in \mathbb{C}^*.$
%\vspace{0.2cm}
%%%%%%%%%%%%%%%%%%%%%%%%%%%%%%%%%%%%%%%%%%%%%%%%%%%%%%%%%%%%%%%%%%%%%%%%%%%
\item[(4)] $\zeta_4: M_kVT_n \longrightarrow \mathrm{GL}_n(\mathbb{C})$ such that
{\small
$$\zeta_4(s_i)=\left( \begin{array}{c|@{}c|c@{}}
   \begin{matrix}
     I_{i-1} 
   \end{matrix} 
      & \textbf{0} & \textbf{0} \\
      \hline
    \textbf{0} &\hspace{0.2cm} \begin{matrix}
   		-1 & 0\\
   		0 & 1\\
   		\end{matrix}  & \textbf{0}  \\
\hline
\textbf{0} & \textbf{0} & I_{n-i-1}
\end{array} \right) \text{ and } \zeta_4(\rho_i^{\alpha})=\left( \begin{array}{c|@{}c|c@{}}
   \begin{matrix}
     I_{i-1} 
   \end{matrix} 
      & \textbf{0} & \textbf{0} \\
      \hline
    \textbf{0} &\hspace{0.2cm} \begin{matrix}
   		0 & \dfrac{1}{y_{\alpha}}\\
   		y_{\alpha} & 0\\
   		\end{matrix}  & \textbf{0}  \\
\hline
\textbf{0} & \textbf{0} & I_{n-i-1}
\end{array} \right)$$
}
$\text{ for } 1\leq i \leq n-1, 0 \leq \alpha \leq k-1, \text{ where } y_{\alpha}\in \mathbb{C}^*.$
%\vspace{0.2cm}
%%%%%%%%%%%%%%%%%%%%%%%%%%%%%%%%%%%%%%%%%%%%%%%%%%%%%%%%%%%%%%%%%%%%%%%%%%%
\item[(5)] $\zeta_5: M_kVT_n \longrightarrow \mathrm{GL}_n(\mathbb{C})$ such that
{\small
$$\zeta_5(s_i)=\left( \begin{array}{c|@{}c|c@{}}
   \begin{matrix}
     I_{i-1} 
   \end{matrix} 
      & \textbf{0} & \textbf{0} \\
      \hline
    \textbf{0} &\hspace{0.2cm} \begin{matrix}
   		-1 & 0\\
   		0 & -1\\
   		\end{matrix}  & \textbf{0}  \\
\hline
\textbf{0} & \textbf{0} & I_{n-i-1}
\end{array} \right) \text{ and } \zeta_5(\rho_i^{\alpha})=\left( \begin{array}{c|@{}c|c@{}}
   \begin{matrix}
     I_{i-1} 
   \end{matrix} 
      & \textbf{0} & \textbf{0} \\
      \hline
    \textbf{0} &\hspace{0.2cm} \begin{matrix}
   		0 & \dfrac{1}{y_{\alpha}}\\
   		y_{\alpha} & 0\\
   		\end{matrix}  & \textbf{0}  \\
\hline
\textbf{0} & \textbf{0} & I_{n-i-1}
\end{array} \right)$$
}
$\text{ for } 1\leq i \leq n-1, 0 \leq \alpha \leq k-1, \text{ where } y_{\alpha}\in \mathbb{C}^*.$
%\vspace{0.2cm}
%%%%%%%%%%%%%%%%%%%%%%%%%%%%%%%%%%%%%%%%%%%%%%%%%%%%%%%%%%%%%%%%%%%%%%%%%%%
\item[(6)] $\zeta_6: M_kVT_n \longrightarrow \mathrm{GL}_n(\mathbb{C})$ such that
{\small
$$\zeta_6(s_i)=\left( \begin{array}{c|@{}c|c@{}}
   \begin{matrix}
     I_{i-1} 
   \end{matrix} 
      & \textbf{0} & \textbf{0} \\
      \hline
    \textbf{0} &\hspace{0.2cm} \begin{matrix}
   		1 & z\\
   		0 & -1\\
   		\end{matrix}  & \textbf{0}  \\
\hline
\textbf{0} & \textbf{0} & I_{n-i-1}
\end{array} \right) \text{ and } \zeta_6(\rho_i^{\alpha})=\left( \begin{array}{c|@{}c|c@{}}
   \begin{matrix}
     I_{i-1} 
   \end{matrix} 
      & \textbf{0} & \textbf{0} \\
      \hline
    \textbf{0} &\hspace{0.2cm} \begin{matrix}
   		0 & \dfrac{1}{y_{\alpha}}\\
   		y_{\alpha} & 0\\
   		\end{matrix}  & \textbf{0}  \\
\hline
\textbf{0} & \textbf{0} & I_{n-i-1}
\end{array} \right)$$
}
$\text{ for } 1\leq i \leq n-1, 0 \leq \alpha \leq k-1, \text{ where } y_{\alpha}, z \in \mathbb{C}^*.$
%\vspace{0.2cm}
%%%%%%%%%%%%%%%%%%%%%%%%%%%%%%%%%%%%%%%%%%%%%%%%%%%%%%%%%%%%%%%%%%%%%%%%%%%
\item[(7)] $\zeta_7: M_kVT_n \longrightarrow \mathrm{GL}_n(\mathbb{C})$ such that
{\small
$$\zeta_7(s_i)=\left( \begin{array}{c|@{}c|c@{}}
   \begin{matrix}
     I_{i-1} 
   \end{matrix} 
      & \textbf{0} & \textbf{0} \\
      \hline
    \textbf{0} &\hspace{0.2cm} \begin{matrix}
   		-1 & z\\
   		0 & 1\\
   		\end{matrix}  & \textbf{0}  \\
\hline
\textbf{0} & \textbf{0} & I_{n-i-1}
\end{array} \right) \text{ and } \zeta_7(\rho_i^{\alpha})=\left( \begin{array}{c|@{}c|c@{}}
   \begin{matrix}
     I_{i-1} 
   \end{matrix} 
      & \textbf{0} & \textbf{0} \\
      \hline
    \textbf{0} &\hspace{0.2cm} \begin{matrix}
   		0 & \dfrac{1}{y_{\alpha}}\\
   		y_{\alpha} & 0\\
   		\end{matrix}  & \textbf{0}  \\
\hline
\textbf{0} & \textbf{0} & I_{n-i-1}
\end{array} \right)$$
}
$\text{ for } 1\leq i \leq n-1, 0 \leq \alpha \leq k-1, \text{ where } y_{\alpha}, z \in \mathbb{C}^*.$
%\vspace{0.2cm}
%%%%%%%%%%%%%%%%%%%%%%%%%%%%%%%%%%%%%%%%%%%%%%%%%%%%%%%%%%%%%%%%%%%%%%%%%%%
\item[(8)] $\zeta_8: M_kVT_n \longrightarrow \mathrm{GL}_{n}(\mathbb{C})$ such that
{\small
\begin{align*}
    \zeta_8(s_i) &= \left( \begin{array}{c|@{}c|c@{}}
   \begin{matrix}
     I_{i-1} 
   \end{matrix} 
      & \textbf{0} & \textbf{0} \\
      \hline
    \textbf{0} &\hspace{0.2cm} \begin{matrix}
   		-a & -\dfrac{(a^2-1)}{b}\\
   		b & a\\
   		\end{matrix}  & \textbf{0}  \\
\hline
\textbf{0} & \textbf{0} & I_{n-i-1}
\end{array} \right) \text{ and },\\
\zeta_8(\rho_i^{\alpha}) &= \left( \begin{array}{c|@{}c|c@{}}
   \begin{matrix}
     I_{i-1} 
   \end{matrix} 
      & \textbf{0} & \textbf{0} \\
      \hline
    \textbf{0} &\hspace{0.2cm} \begin{matrix}
   		0 & \dfrac{1}{y_{\alpha}}\\
   		y_{\alpha} & 0\\
   		\end{matrix}  & \textbf{0}  \\
\hline
\textbf{0} & \textbf{0} & I_{n-i-1}
\end{array} \right)
\end{align*}
% $$\zeta_8(s_i)=\left( \begin{array}{c|@{}c|c@{}}
%    \begin{matrix}
%      I_{i-1} 
%    \end{matrix} 
%       & \textbf{0} & \textbf{0} \\
%       \hline
%     \textbf{0} &\hspace{0.2cm} \begin{matrix}
%    		-a & -\dfrac{(a^2-1)}{b}\\
%    		b & a\\
%    		\end{matrix}  & \textbf{0}  \\
% \hline
% \textbf{0} & \textbf{0} & I_{n-i-1}
% \end{array} \right) \text{ and } $$
% $$\zeta_8(\rho_i^{\alpha})=\left( \begin{array}{c|@{}c|c@{}}
%    \begin{matrix}
%      I_{i-1} 
%    \end{matrix} 
%       & \textbf{0} & \textbf{0} \\
%       \hline
%     \textbf{0} &\hspace{0.2cm} \begin{matrix}
%    		0 & \dfrac{1}{y_{\alpha}}\\
%    		y_{\alpha} & 0\\
%    		\end{matrix}  & \textbf{0}  \\
% \hline
% \textbf{0} & \textbf{0} & I_{n-i-1}
% \end{array} \right)$$
}
$\text{ for } 1\leq i \leq n-1, 0 \leq \alpha \leq k-1, \text{ where } a\in \mathbb{C}$ and $b, y_{\alpha}\in \mathbb{C}^*.$
\end{itemize}

% $$\zeta_1(\rho_i^1)=\left( \begin{array}{c|@{}c|c@{}}
%    \begin{matrix}
%      I_{j-1} 
%    \end{matrix} 
%       & \textbf{0} & \textbf{0} \\
%       \hline
%     \textbf{0} &\hspace{0.2cm} \begin{matrix}
%    		0 & \dfrac{1}{a_1}\\
%    		a_1 & 0\\
%    		\end{matrix}  & \textbf{0}  \\
% \hline
% \textbf{0} & \textbf{0} & I_{n-j}
% \end{array} \right), \dots, 
% \zeta_1(\rho_i^{k-1})=\left( \begin{array}{c|@{}c|c@{}}
%    \begin{matrix}
%      I_{j-1} 
%    \end{matrix} 
%       & \textbf{0} & \textbf{0} \\
%       \hline
%     \textbf{0} &\hspace{0.2cm} \begin{matrix}
%    		0 & \dfrac{1}{a_{k-1}}\\
%    		a_{k-1} & 0\\
%    		\end{matrix}  & \textbf{0}  \\
% \hline
% \textbf{0} & \textbf{0} & I_{n-j}
% \end{array} \right),$$
\end{theorem}

\begin{theorem}
Let $
\zeta : M_kVT_n \longrightarrow \mathrm{GL}_n(\mathbb{C})$
be a complex homogeneous \(2\)-local representation of \(M_kVT_n\) with
$k\geq 1$ and $n \ge 3$. By Theorem~\ref{thm 41}, the representation $\zeta$ is equivalent to one of the eight
representations $\zeta_i$, $1 \le i \le 8$, described therein.
The reducibility of $\zeta$ is characterized as follows.
\begin{itemize}
    \item[(a)] If $\zeta$ is equivalent to $\zeta_1$, then $\zeta$ is reducible.
    \item[(b)] If $\zeta$ is equivalent to one of $\zeta_2, \zeta_3, \zeta_4,$ or $\zeta_5$,
    then $\zeta$ is reducible if and only if
    $y_{\beta_1}=y_{\beta_2}$ for all $0 \le \beta_1, \beta_2 \le k-1$.
    \item[(c)] If $\zeta$ is equivalent to one of $\zeta_6$ or $\zeta_7$,
    then $\zeta$ is irreducible.
    \item[(d)] If $\zeta$ is equivalent to $\zeta_8$, then $\zeta$ is reducible
    if and only if $y_{\beta_1}=y_{\beta_2}:=y$ for all $0 \le \beta_1, \beta_2 \le k-1$ and $\left( \dfrac{b}{y}=1+a \text{ \ or\ \  } \dfrac{b}{y}=1-a\right).$
\end{itemize}
\end{theorem}

\begin{proof}
We analyze each case separately, omitting the case (a) as it is trivial.
\begin{itemize}
\item[(b)]
Assume that $\zeta$ is equivalent to one of
$\zeta_2, \zeta_3, \zeta_4,$ or $\zeta_5$.
It suffices to treat the case $\zeta \simeq \zeta_2$,
as the remaining cases follow by similar arguments.

For the necessary condition, 
suppose that there exist $0 \le \beta_1 \neq \beta_2 \le k-1$ such that
$y_{\beta_1} \neq y_{\beta_2}$.
Let $\zeta_2'$ be the representation equivalent to $\zeta_2$ defined by
\[
\zeta_2'=P^{-1}\zeta_2P,
\qquad
P=\mathrm{diag}\big(y_{\beta_1}^{1-n},\,y_{\beta_1}^{2-n},\ldots,y_{\beta_1}^{-1},1\big).
\]
A direct computation shows that, for the generators $\rho_i^{\beta_1}$ and $\rho_i^{\beta_2}$ with $1\le i\le n-1$, the images under $\zeta_2'$ are given by $$\zeta_2'(\rho_i^{\beta_1})=\left( \begin{array}{c|@{}c|c@{}} \begin{matrix} I_{i-1} \end{matrix} & \textbf{0} & \textbf{0} \\ \hline \textbf{0} &\hspace{0.2cm} \begin{matrix} 0 & 1\\ 1 & 0\\ \end{matrix} & \textbf{0} \\ \hline \textbf{0} & \textbf{0} & I_{n-i-1} \end{array} \right)$$
and
$$\zeta_2'(\rho_i^{\beta_2})=\left( \begin{array}{c|@{}c|c@{}} \begin{matrix} I_{i-1} \end{matrix} & \textbf{0} & \textbf{0} \\ \hline \textbf{0} &\hspace{0.2cm} \begin{matrix} 0 & \dfrac{y_{\beta_1}}{y_{\beta_2}}\\ \dfrac{y_{\beta_2}}{y_{\beta_1}} & 0\\ \end{matrix} & \textbf{0} \\ \hline \textbf{0} & \textbf{0} & I_{n-i-1} \end{array} \right).$$
Assume, toward a contradiction, that $\zeta_2$ is reducible.
Then $\zeta_2'$ is reducible and so it admits a nontrivial proper invariant subspace
$U \subset \mathbb{C}^n$.
Let $u=(u_1,u_2,\ldots,u_n)^T \in U$ be nonzero. For each $1 \le i \le n-1$, we have
\[
\zeta_2'(\rho_i^{\beta_1})u - u
= (u_{i+1}-u_i)(e_i-e_{i+1}) \in U.
\]
If $u_i=u_{i+1}$ for all $1\leq i\leq n-1$, then $u$ is a scalar multiple of
$(1,1,\ldots,1)^T$, which is impossible since
$U$ is invariant under $\zeta_2'(\rho_i^{\beta_2})$ and
$y_{\beta_1} \neq y_{\beta_2}$.
Hence, there exists $1 \le j \le n-1$ such that $u_j\neq u_{j+1}$, and so we get that $e_j-e_{j+1} \in U$.
By applying $\zeta_2'(\rho_{j+1}^{\beta_1})$
and $\zeta_2'(\rho_{j-1}^{\beta_1})$ to $e_j-e_{j+1}$, we obtain
$$
\zeta_2'(\rho_{j+1}^{\beta_1})(e_j - e_{j+1})-(e_j - e_{j+1}) = e_{j+1} - e_{j+2} \in U
$$ 
and
$$\zeta_2'(\rho_{j-1}^{\beta_1})(e_j - e_{j+1})-(e_j - e_{j+1}) = e_{j-1} - e_j \in U.
$$
Continuing in the process we get 
\[
e_i-e_{i+1} \in U \quad \text{for all } 1 \le i \le n-1.
\]
Thus, no standard basis vector $e_i$ lies in $U$; otherwise,
$U=\mathbb{C}^n$, a contradiction.
Next,
\[
\zeta_2'(\rho_1^{\beta_2})(e_1-e_2)
+ \dfrac{y_{\beta_2}}{y_{\beta_1}}(e_1-e_2)
= \left(\dfrac{y_{\beta_2}}{y_{\beta_1}}
- \dfrac{y_{\beta_1}}{y_{\beta_2}}\right)e_1 \in U.
\]
Since $e_1 \notin U$, it follows that
$y_{\beta_1}=\pm y_{\beta_2}$.
As $y_{\beta_1}\neq y_{\beta_2}$, we must have
$y_{\beta_1}=-y_{\beta_2}$.
In this case we have,
\[
\zeta_2'(\rho_1^{\beta_1})(e_2-e_3)
-\zeta_2'(\rho_1^{\beta_2})(e_2-e_3)=2e_1 \in U,
\]
which implies a contradiction since $e_1\notin U$. Thus, $\zeta_2$ is irreducible whenever
$y_{\beta_1}\neq y_{\beta_2}$ for some $0\leq \beta_1 \neq \beta_2\leq k$, as required.\vspace{0.1cm}

For the sufficient condition, suppose that we have $y_{\beta_1}=y_{\beta_2}$ for all $0\leq \beta_1,\beta_2\leq k-1$. Then the vector $(1,1,\ldots,1)^T$ is invariant under the action of $\zeta_2'(s_i)$ and $\zeta_2'(\rho_i^\alpha)$ for all $1\leq i \leq n$ and \(0 \le \alpha \le k-1\). Hence, $\zeta_2'$ admits a nontrivial invariant subspace and is therefore reducible. Thus, $\zeta_2$ is also reducible. \vspace{0.1cm}
\item[(c)]
Assume that $\zeta$ is equivalent to $\zeta_6$ or $\zeta_7$.
Without loss of generality, suppose that $\zeta \simeq \zeta_6$.
We distinguish the following two cases.
\begin{itemize}
\item[(i)]
There exist $0\leq \beta_1 \neq \beta_2 \leq k-1$ such that $y_{\beta_1} \neq y_{\beta_2}$.
In this situation, repeating the argument of case (b) shows that
$\zeta$ is irreducible.
\item[(ii)]
$y_{\beta_1} = y_{\beta_2} := y$ for all $0 \le \beta_1, \beta_2 \le k-1$.
In this case, we consider an equivalent representation defined as follows.
Let $\zeta_6'$ be the representation equivalent to $\zeta_6$ given by
\[
\zeta_6'=Q^{-1}\zeta_6Q,
\qquad
Q=\mathrm{diag}\big(y^{1-n},\,y^{2-n},\ldots,y^{-1},1\big).
\]
A direct computation shows that, for the generators $s_i$ and $\rho_i^\alpha$
with $1 \le i \le n-1$ and $0 \le \alpha \le k-1$, we have
\[
\zeta_6'(s_i)=
\left(
\begin{array}{c|@{}c|c@{}}
\begin{matrix} I_{i-1} \end{matrix} & \mathbf{0} & \mathbf{0} \\
\hline
\mathbf{0} &
\begin{matrix}
\ 1 & yz\\
\ 0 & -1
\end{matrix}
& \mathbf{0} \\
\hline
\mathbf{0} & \mathbf{0} & I_{n-i-1}
\end{array}
\right),
\]
and
\[
\zeta_6'(\rho_i^{\alpha})=
\left(
\begin{array}{c|@{}c|c@{}}
\begin{matrix} I_{i-1} \end{matrix} & \mathbf{0} & \mathbf{0} \\
\hline
\mathbf{0} &
\begin{matrix}
\ 0 & 1\\
\ 1 & 0
\end{matrix}
& \mathbf{0} \\
\hline
\mathbf{0} & \mathbf{0} & I_{n-i-1}
\end{array}
\right).
\]
Assume, toward a contradiction, that $\zeta_6$ is reducible.
Then $\zeta_6'$ is also reducible and so it admits a non-trivial invariant subspace
$U \subset \mathbb{C}^n$.
As in case (b), invariance under
$\zeta_6'(\rho_i^{\alpha})$, $0 \le \alpha \le k-1$,
implies that $e_i - e_{i+1} \in U$ for all $1 \le i \le n-1$.
Consequently, no standard basis vector $e_i$ belongs to $U$;
otherwise, $U = \mathbb{C}^n$, a contradiction. On the other hand, we have
\[
\zeta_6'(s_1)(e_2 - e_3) + (e_2 - e_3) = yz\, e_1,
\]
with $y \neq 0$ and $z \neq 0$, which implies that $e_1 \in U$.
Hence $U = \mathbb{C}^n$, a contradiction.
Therefore, $\zeta_6$ is irreducible.
\end{itemize}

\item[(d)] Suppose now that $\zeta$ is equivalent to $\zeta_8$.

For the necessary condition, assume that $\zeta_8$ is reducible. A work similar to that in case (b) implies that $y_{\beta_1} = y_{\beta_2} := y$ for all $0 \le \beta_1, \beta_2 \le k-1$. So, we consider an equivalent representation defined in a similar way as in case (c). Let $\zeta_8'$ be the representation equivalent to $\zeta_8$ given by
\[
\zeta_8'=Q^{-1}\zeta_8Q,
\qquad
Q=\mathrm{diag}\big(y^{1-n},\,y^{2-n},\ldots,y^{-1},1\big).
\]
A direct computation shows that, for the generators $s_i$ and $\rho_i^\alpha$
with $1 \le i \le n-1$ and $0 \le \alpha \le k-1$, the actions under $\zeta_8'$ are given as follows.
\[
\zeta_8'(s_i)=
\left(
\begin{array}{c|@{}c|c@{}}
\begin{matrix} I_{i-1} \end{matrix} & \mathbf{0} & \mathbf{0} \\
\hline
\mathbf{0} &
\begin{matrix}
\ -a & - \dfrac{(a^2-1)y}{b}\\
\ \dfrac{b}{y} & a
\end{matrix}
& \mathbf{0} \\
\hline
\mathbf{0} & \mathbf{0} & I_{n-i-1}
\end{array}
\right),
\]
and
\[
\zeta_8'(\rho_i^{\alpha})=
\left(
\begin{array}{c|@{}c|c@{}}
\begin{matrix} I_{i-1} \end{matrix} & \mathbf{0} & \mathbf{0} \\
\hline
\mathbf{0} &
\begin{matrix}
\ 0 & 1\\
\ 1 & 0
\end{matrix}
& \mathbf{0} \\
\hline
\mathbf{0} & \mathbf{0} & I_{n-i-1}
\end{array}
\right).
\]
As $\zeta_8$ is reducible, we have $\zeta'_8$ is also reducible, and so it admits a nontrivial invariant subspace
$U \subset \mathbb{C}^n$. Let $u=(u_1,u_2,\ldots,u_n)^T \in U$ be nonzero. For each $1 \le i \le n-1$, we have
\[
\zeta_2'(\rho_i^{\alpha})u - u
= (u_{i+1}-u_i)(e_i-e_{i+1}) \in U.
\]
We now consider two cases as follows.
\begin{itemize}
    \item[(i)] Suppose that there exists $1 \le j \le n-1$ such that that $u_j\neq u_{j+1}$, then we get, as in case (b), that
\[
e_i-e_{i+1} \in U \quad \text{for all } 1 \le i \le n-1.
\]
In this case, we can see that 
$$\zeta_8'(s_1)(e_1-e_2)+\zeta_8'(s_1)(e_2-e_3)-(e_2-e_3)+a(e_1-e_2)=\left( \dfrac{b}{y}-1-a \right)e_2\in U.$$
But $e_2$ can not be in $U$, otherwise we get $U=\mathbb{C}^n$, which is impossible. Hence, we get that $ \dfrac{b}{y}-1-a=0$, and so $ \dfrac{b}{y}=1+a$.
\item[(ii)] Suppose now $u_j= u_{j+1}$ for all $1\leq j\leq n-1$. Then $U$ is generated by the vector $(1,1,\ldots ,1)^T$. Similarly, as in (i), we get that $\dfrac{b}{y}=1-a.$
\end{itemize}
Thus, we get that either $ \dfrac{b}{y}=1+a$ or $\dfrac{b}{y}=1-a,$ as required.

For the sufficient condition, suppose that $y_{\beta_1}=y_{\beta_2}$ for all $0 \le \beta_1, \beta_2 \le k-1$. We consider two cases as follows.
\begin{itemize}
    \item[(i)] If $\dfrac{b}{y}=1+a$, then one readily verifies that the vector $(1,1,\ldots,1)^T$ is invariant on the matrices $\zeta_8'(s_i)$ and $\zeta_8'(\rho_i^{\alpha})$ by left multiplication, for all $1\leq i\leq n-1$ and $0\leq \alpha \leq k-1$.
    
    \item[(ii)] If $\dfrac{b}{y}=1-a$, then one readily verifies that the vector $(1,1,\ldots,1)^T$ is invariant on the matrices $\zeta_8'(s_i)$ and $\zeta_8'(\rho_i^{\alpha})$ by right multiplication, for all $1\leq i\leq n-1$ and $0\leq \alpha \leq k-1$.
\end{itemize}

Hence, $\zeta_8'$ admits a nontrivial invariant subspace in both cases and is therefore reducible. Thus, $\zeta_8$ is also reducible.
\end{itemize}
\end{proof}

\end{document}
